\chapter{Ethereum Virtual Machine (EVM)}
Eterijum \cite{wood2014ethereum} \eng{Ethereum} predstavlja prvu blokčejn platformu koja je uvela Tjuring kompletne pametne ugovore i kao takva se u savremenoj akademskoj literaturi i industrijskoj praksi prepoznaje kao jedna od najaktivnijih i tehnološki najrazvijenijih platformi za razvoj decentralizovanih aplikacija \cite{hashkey2024}. Empirijske analize aktivnosti programerske zajednice ukazuju na izuzetno visok intenzitet razvoja, koji se ogleda u velikom broju izmena \eng{commit} u ključnim repozitorijumima povezanih sa Ethereum ekosistemom. Ovi pokazatelji potvrđuju centralnu ulogu Ethereuma u procesu tehnoloških inovacija unutar šireg blokčejn prostora \cite{vaccargiu2025mining}.

Počevši od 15. septembra 2022. godine, mreža Eterijuma koristi mehanizam konsenzusa zasnovan na dokazivanju uloga \eng{Proof of Stake, PoS}, u okviru koga validatori imaju ključnu ulogu u procesu predlaganja i validacije novih blokova. Da bi učesnik mreže stekao status validatora, neophodno je da zaključa najmanje 32 ETH, koji predstavljaju primarnu kriptovalutu Eterijum blokčejna. Prema dostupnim podacima, od drugog kvartala 2024. godine mreža obuhvata približno jedan milion aktivnih validatora, uz umerene fluktuacije tog broja tokom vremena \cite{ethereumvalidators}.

Analiza ponašanja validatora u dužem vremenskom periodu \cite{grandjean2024ethereum} ukazuje na izuzetno visok stepen učešća u konsenzusnom procesu, sa prosečnim nivoom participacije od oko 99\%. Povremeni padovi u učešću uglavnom su povezani sa planiranim nadogradnjama protokola ili tehničkim problemima u softverskim klijentima koje validatori koriste. Pored toga, ista studija identifikuje visok nivo decentralizacije konsenzusnog sloja mreže, kvantifikovan korišćenjem standardnih metrika kao što su Nakamotov koeficijent, Ginijev koeficijent i Herfindahl-Hiršmanov indeks, čime se dodatno potvrđuje robusnost i otpornost Ethereum mreže na koncentraciju moći među učesnicima.

Eterijum se najčešće modeluje kao sistem sastavljen od dva protokolarno razdvojena sloja, koji funkcionišu kao međusobno koordinisane, ali logički nezavisne mreže ravnopravnih čvorova. Svaki sloj održava sopstveno lokalno stanje relevantno za svoju funkciju u okviru protokola:
\begin{enumerate}
\item \textbf{Konsenzusni sloj} \eng{Consensus Layer}, poznat i kao \textit{lanac svetionika} \eng{Beacon Chain}, definiše pravila za izbor predlagača blokova, određivanje kanonske istorije lanca, kao i mehanizme glasanja i finalnosti blokova.
\item \textbf{Izvršni sloj} \eng{Execution Layer} predstavlja globalni distribuirani računarski sloj sistema i zadužen je za izvršavanje transakcija, kao i za definisanje pravila interakcije i ažuriranja stanja Etereum virtuelne mašine.
\end{enumerate}

Svaki validator mora imati funkcionalni pristup i konsenzusnom i izvršnom sloju protokola. Konsenzusni sloj se uvek pokreće lokalno, dok se izvršni sloj može realizovati kao lokalni klijent ili kao udaljeni eksterni servis povezan putem standardizovanog interfejsa. U praksi, validatori najčešće pokreću oba sloja lokalno, budući da oslanjanje na udaljeni izvršni klijent povećava napadnu površinu i uvodi dodatne bezbednosne rizike.

Komunikacija između konsenzusnog i izvršnog sloja je dvosmerna i ostvaruje se isključivo putem standardizovanog interfejsa poznatog kao \textit{Engine API}, koji predstavlja jedini dozvoljeni mehanizam njihove koordinacije. Pod pretpostavkom odsustva vizantijskog ponašanja validatora, proces predlaganja i validacije novog bloka odvija se kroz sledeće faze:
\begin{enumerate}
\item \label{PoS-1} Validator izabran za predlagača bloka u okviru konsenzusnog sloja inicira zahtev ka povezanom izvršnom klijentu radi generisanja sadržaja novog bloka.
\item \label{PoS-2} Izvršni klijent bira nepotvrđene transakcije iz svog lokalnog skupa, deterministički ih izvršava u okviru lokalne Eterijum virtuelne mašine i konstruiše izvršni sadržaj bloka \eng{Execution Payload}.
\item \label{PoS-3} Konsenzusni sloj predlagača preuzima izvršni sadržaj i proverava njegovu strukturnu, formatnu i kriptografsku ispravnost u skladu sa važećim pravilima protokola.
\item \label{PoS-4} Predlagač bloka kombinuje izvršni sadržaj sa konsenzusnim podacima u jedinstvenu strukturu bloka, koju zatim kriptografski potpisuje i distribuira svim validatorima u konsenzusnom sloju.
\item \label{PoS-5} Ostali validatori konsenzusnog sloja primaju predloženi blok i verifikuju isključivo konsenzusni deo bloka u skladu sa pravilima protokola.
\item \label{PoS-6} Validatori emitovanjem potvrda učestvuju u procesu potvrđivanja bloka i njegovog uključivanja u kanonsku istoriju blokčejna.
\item \label{PoS-7} Nakon odluke konsenzusnog sloja o prihvatanju bloka, izvršni sloj svakog validatora deterministički ponovo izvršava transakcije sadržane u izvršnom sadržaju bloka i ažurira svoje lokalno stanje sistema.
\end{enumerate}

\section{Konsenzusni sloj}
Vreme u konsenzusnom sloju podeljeno je na epohe \eng{epochs}, pri čemu svaka epoha traje 6 minuta i 24 sekunde i sastoji se od 32 slota, dok trajanje pojedinačnog slota iznosi 12 sekundi. Na početku svake epohe, svi aktivni validatori se pseudonasumično raspoređuju u odbore \eng{committees} približno jednake veličine, tako da svakom slotu odgovara tačno jedan odbor. Radi očuvanja bezbednosnih garancija protokola, veličina svakog odbora ne može biti manja od 128 validatora.

U svakom slotu se pseudonasumično bira jedan validator (Faza \ref{PoS-1}) koji preuzima ulogu predlagača bloka. Predlagač je odgovoran za kreiranje novog bloka tako što, putem povezanog izvršnog klijenta, bira transakcije iz svog lokalnog skupa nepotvrđenih transakcija i uključuje ih u izvršni sadržaj bloka (Faza \ref{PoS-2}). Nakon interne provere izvršnog sadržaja (Faza \ref{PoS-3}) i njegove integracije sa konsenzusnim podacima u jedinstvenu strukturu bloka (Faza \ref{PoS-4}), predloženi blok se distribuira ostalim validatorima u odgovarajućem odboru (Faza \ref{PoS-5}).

Preostali validatori u odboru deluju kao glasači (Faza \ref{PoS-6}) i u svakom slotu mogu emitovati najviše jednu potvrdu \eng{attestation}, kojom potvrđuju ispravnost konsenzusnog dela predloženog bloka. U ovoj fazi se ne proverava validnost transakcija niti izvršni sadržaj bloka, već isključivo usklađenost sa pravilima konsenzusnog protokola. Kako se svaki odbor sastoji od velikog broja validatora, direktna propagacija pojedinačnih potvrda bi dovela do zagušenja mreže, u skladu sa Teoremom \ref{teorema:gunter}. Iz tog razloga, svaki odbor je dodatno podeljen na manje komunikacione podmreže \eng{subnets}, pri čemu se svaki validator nalazi u tačno jednoj podmreži.

Unutar svake podmreže, šesnaest validatora se bira u ulogu agregatora \eng{aggregators}, koji prikupljaju potvrde ostalih validatora i kombinuju njihove kriptografske potpise u jedan agregatni potpis primenom BLS šeme agregacije potpisa i emituju dobijeni agregatni potpis široj mreži. Ovakav mehanizam omogućava efikasnu verifikaciju velikog broja potvrda uz minimalan komunikacioni i računarski trošak.

Validator koji je izabran za predlagača bloka uključuje u predloženi blok agregirane BLS potvrde iz svake komunikacione podmreže, pri čemu se za svaku podmrežu bira agregacija sa najvećom podrškom validatora. Ove agregirane potvrde se kombinuju u jedan konačni BLS potpis, koji postaje sastavni deo strukture bloka.

Međutim, usled asinhrone prirode distribuirane mreže, predloženi blok ne dospeva do svih validatora istovremeno. Kao posledica toga, pojedini validatori mogu primiti novi blok tek u narednom slotu. Validatori koji nisu primili novi blok nastavljaju da glasaju za poslednji blok koji im je poznat. Ukoliko novi predlagač nije upoznat sa prethodno predloženim blokom, on može konstruisati blok koji se nadovezuje direktno na poslednji finalizovani blok, čime dolazi do privremenog račvanja lanca.

Pretpostavimo da je poslednji finalizovani blok označen sa $B$. Predlagač $P$ predlaže blok $B_1$, koji ne dospeva do predlagača $Q$ u narednom slotu. Predlagač $Q$ zatim predlaže blok $B_2$ koji se takođe nadovezuje na blok $B$, čime nastaje račvanje. U narednim slotovima, naredni predlagači mogu nastaviti izgradnju lanca na obe grane, pri čemu se, na primer, tri nova bloka nadovezuju na granu $B_1$, dok se dva bloka linearno nadovezuju na granu $B_2$ (Slika TODO).

S obzirom na to da svaka potvrda sadrži informaciju o slotu u kome je emitovana, validator u svakom narednom slotu ima mogućnost da promeni svoj izbor i emituje novu potvrdu za nefinalizovani blok na suprotnoj grani, ako ne želi da glasa za istu granu. Iz tog razloga, u procesu izbora glavne grane u obzir se uzima isključivo poslednja važeća potvrda svakog validatora, odnosno potvrda sa najvećim brojem slota, koja je uključena u agregirane BLS potpise. Ovaj princip poznat je kao \eng{Latest Message Driven, LMD}.

Kako se u Eterijum mreži zaključava primarna kriptovaluta u vidu uloga validatora, izbor grane koja će biti proglašena glavnom u slučaju račvanja zasniva se na \textit{pravilu najtežeg lanca}, koje predstavlja varijaciju klasičnog pravila najdužeg lanca. Ovo pravilo ne posmatra broj blokova na grani, već ukupnu količinu ETH zaključanu u podstablu posmatranog bloka, uzimajući u obzir poslednje potvrde svih validatora koji su glasali za njega ili njegove potomke \eng{Greedy Heaviest-Observed Sub-Tree, GHOST}.

U slučaju račvanja, Eterijum blokčejn obezbeđuje verovatnosnu finalnost primenom kombinovanog mehanizma \textit{LMD-GHOST}, koji koristi poslednje potvrde validatora za deterministički izbor grane sa najvećom ukupnom težinom.